\chapter{解题感悟}
\label{cp:notes}

这周的练习和上几周相比，感觉重点变了：前几周学的是一堆单独的命令，知道 \verb|cd|、\verb|cp|、\verb|chmod| 各自怎么用；这周的四道题更像是要把这些命令串成一条条流程——先从复制工程、清理缓存开始，再把检查动作固化成 \verb|check.sh|，然后用 Makefile 描述产物之间的依赖，最后在虚拟环境里把 wheel 构建出来。真正动手之后才意识到，一条命令会敲和一条流程能稳定跑通，中间差的是不少细节。

让我印象最深的是 Make 那道题。以前一直以为 \verb|make| 就是把几条命令打包起来按顺序执行，直到看见第二次运行它只回了一句 \texttt{Nothing to be done for `all'.}，而 \verb|touch data.csv| 之后它又只重跑了受影响的两步，才明白它判断的其实是「依赖文件是不是比目标新」，而不是「我是不是刚敲过这条命令」。把依赖关系写完整之后，工具就能替我做增量判断；反过来，如果依赖漏写了，它就会在需要重跑的时候偷懒，这种错误还特别不容易被发现。

前几周配好的虚拟环境，这周才算真正派上用场。\verb|make build| 第一次执行时直接报 \texttt{No module named build}，当时有点懵，后来在 \verb|(.venv)| 环境里再跑一次就顺利产出了 wheel。这件事让我理解了第一周讲的「依赖必须隔离」并不是形式要求：构建工具也是依赖，它装在哪，决定了这条构建命令在别人的机器上还能不能跑得起来。同样让我有收获的是 \verb|check.sh| 里的 \verb|set -euo pipefail| 和最后那个 \verb|echo $?|——门禁的价值在于把「通过还是没通过」压缩成一个退出码，而不是靠人盯着满屏输出自己判断。

具体操作上，这周也踩了一些小坑。jq 的表达式一开始完全靠试，写过滤条件时把 \verb|select| 和管道的位置搞混过好几次，最后是照着讲义里「先展开、再过滤、最后排序」的顺序一步步写才顺下来；vim 里改文件时也还不太熟练，经常要先 \verb|cat| 一遍确认真的改对了。好在这周的截图我是边做边存的，写报告时不用担心漏掉某一步的操作记录，这一点比之前从容很多。

写报告本身也顺利了一些。上周刚学会用 LaTeX 写东西的时候，光是调格式就花了不少时间；这次图片、代码块和交叉引用基本都是照搬上次的写法，用 \verb|\label| 和 \verb|\ref| 引用图片也比手动对着编号数要省心，所以能把更多时间花在把每一步的实验结果说清楚上。当然不足还是很明显的：查资料的时间依然不少，尤其是 jq 和 Make 的语法细节；有些快捷键和环境配置也只是照做，还没到能给别人解释清楚的程度。希望接下来几周能把这些常用操作练到不用查资料，让实验的重心从「跑通命令」慢慢转到「想清楚流程」上。
